\documentclass[11pt]{article}

\usepackage{maiacustom}

\begin{document}

\psettitle{Astronomy Question Bank}

These questions were carefully written or selected to prepare students for the astronomy team selection process in Brazil. Some questions are not original; those are marked with () before the statement. The question-bank template is the same one used by Professor Kevin Zhou. His work is invaluable, and many ideas in this set can be found in his handouts.

% \begin{psidea}{Idea title}{}
% Idea
% \end{psidea}

% \begin{psexample}{Example title}{}
% Example
% \end{psexample}

% \begin{pssolution*}{}{}
% Solution
% \end{pssolution*}

% \begin{psremark*}{Remark title}{}
% Remark
% \end{psremark*}
\section{Photometry and Modern Physics}

\pts{3} 
\begin{pproblem}
    The goal of this question is to derive an expression for the \textit{optical depth}. Imagine that a beam of light passes through a region of space containing a certain number of particles. If \(I_0\) is the intensity of the light before it passes through that region of space and \(I\) is the intensity afterward, the optical depth is defined as:
    \[I = I_0e^{-\tau}\]
    where \(\tau\) is the optical depth.
    \\
    Consider a region of space with particle number density \(n\).
    \begin{alternativas}
        \item What is the number of particles in an area \(A\) and thickness \(dz\)?
        \item Assuming that each particle has cross section \(\sigma\), what is the area blocked by the particles?
        \item Find the formula for \(\tau\). 
    \end{alternativas}

\begin{pssolution*}{}{}
    \begin{alternativas}
        \item The volume number density of particles is \(n\), so the number \(dN\) of particles in a volume \(dV\) is
        \[\boxed{dN = ndV \equiv nAdz}\]

        \item If each particle occupies an area \(\sigma\), the area occupied by \(dN\) particles is
        \[\boxed{dS = \sigma dN = \sigma nAdz}\]

        \item The area occupied by the particles is expected to obstruct the passage of light. Since the intensity is proportional to the area available for the light to pass through, we have
        \[\frac{dI}{I} = -\frac{dS}{A} = -n\sigma dz\]
        Integrating, we obtain

        \[I = I_0 e^{-n\sigma z}\]

        That is,

        \[\boxed{\tau = n\sigma z}\]
    \end{alternativas}
\end{pssolution*}
\end{pproblem}

\pts{2}
\begin{pproblem} The Triangulum Galaxy, \(M33\), is the third-largest galaxy in the Local Group. It lies at a distance \(d=970 \text{ kpc}\) from us and has an apparent magnitude of \(5.72\). Given that it contains approximately \(40\) billion stars, find the average luminosity of the stars in \(M33\). Does your estimate seem consistent with reality? Why?
\begin{pssolution*}{}{}
    Using Pogson's equation

    \[m-M_\odot = -2.5\log \left(\frac{L_{g}}{L_\odot}\left(\frac{10 \text{ pc}}{d}\right)^2\right)\]

    Solving for \(L_g\),

    \[L_g = L_\odot \left(\frac{d}{10\text{ pc}}\right)^2 10^{-0.4(m-M_\odot)}\]

    We also have \(L_g = \overline{L_\star} N\), and thus

    \[\overline{L_\star} = \frac{L_\odot}{N} \left(\frac{d}{10\text{ pc}}\right)^2 10^{-0.4(m-M_\odot)}\]

    Using the values from the table of constants, we obtain

    \[\overline{L_\star} \approx 0.1L_\odot\]

    This estimate does not match reality, since factors such as interstellar extinction increase the apparent magnitude of \(M33\). A consistent estimate would be of the same order of magnitude as the luminosity of the Sun.

\end{pssolution*}
\end{pproblem}

\pts{4}
\begin{pproblem}
    Consider a universe with a constant stellar number density, that is, a number of stars per unit volume, of value \(n\). Assume that every star has \(L = L_\odot\) and that the universe contains a total of \(N_0\) stars.
    
    \begin{alternativas}
    \item What is the probability that the absolute magnitude of a star, seen from the center of the universe, lies between \(m\) and \(m+dm\), where \(dm\) is an infinitesimal interval of magnitude? Leave your answer in terms of \(m\) and the largest possible magnitude, \(m_{lim}\), of a star at the "edge" of the universe.

    \item What is the probability that a star can be observed with the naked eye?
    
    \textbf{Data:} \[\int_{-\infty}^{6}10^{0.6(x-A)}dx \approx 2881.6 e^{-1.381 A}\]
    \end{alternativas}
\begin{pssolution*}{}{}
    \begin{alternativas}
    \item The apparent magnitude of a star is given by

    \[m = -2.5\log F + C\]

    where \(C\) is a constant. Writing \(F = L/4\pi r^2\), we have

    \[m = -2.5\log\left(\frac{L}{4\pi r^2}\right) + C = 5\log r -2.5 \log \left(\frac{L}{4\pi}\right)+ C\]

    Since \(L\) is the same for every star, we can absorb the second term into the constant. Thus

    \[m=5\log r + C \rightarrow r = 10^{0.2(m-C)}\]

    Differentiating the expression for \(m\), we have

    \[\frac{dm}{dr} = \frac{5}{r\ln 10}\]

    Now we find the probability that a star lies at a distance \(r\).

    The number of stars between \(r\) and \(r+dr\) is

    \[dN(r) = 4\pi n r^2dr\]

    Dividing by \(N_0\), the probability that a star lies at that distance is

    \[dP(r) = \frac{dN(r)}{N_0} = \frac{4\pi n r^2}{N_0}dr \rightarrow dr = \frac{N_0}{4\pi n r^2} dP(r)\]

    Substituting the expression for \(dr\),

    \[\frac{r\ln 10}{5}dm = \frac{N_0}{4\pi n r^2}dP(r)\]

    Isolating \(dP(r)\),

    \[dP(r) = \frac{4\pi \ln 10 n r^3}{5 N_0}dm\]

    Substituting the expression for \(r\), we change variables \(r\rightarrow m\) in \(dP\):

    \[dP(m) = \frac{4\pi n \ln 10  }{5 N_0}10^{0.6(m-C)}dm\]

    To find \(C\), compare with the limiting magnitude at the "edge" of the universe. Since the universe contains \(N_0\) stars, its radius is given by

    \[N_0 = \frac{4\pi R^3}{3} \rightarrow R = \left(\frac{3N_0}{4\pi}\right)^{1/3}\]

    Writing the equation for \(m_{lim}\),
    
    \[m_{lim} = 5\log R + C\]

    we can isolate the constant \(C = m_{lim}-5\log R\). Substituting into the expression for \(dP(m)\),

    \[dP(m) = \frac{4\pi n \ln 10  }{5 N_0}10^{0.6(m-m_{lim}+5\log R)}dm\]

    Working through this expression,

    \[dP(m) = \frac{4\pi n \ln 10  }{5 N_0}10^{0.6(m-m_{lim})}10^{3\log R}dm\]

    Using the logarithm identities \(a\log b = \log b^a\) and \(10^{\log x} = x\), we have

    \[dP(m) = \frac{4\pi n R^3\ln 10}{5N_0}10^{0.6(m-m_{lim})}dm\]

    Substituting \(R\),

    \[\boxed{dP(m) = \frac{3\ln 10}{5}10^{0.6(m-m_{lim})}dm}\]

    The probability therefore depends on neither \(n\) nor \(N_0\).

    \item For a star to be visible, its magnitude must satisfy \(m\leq 6\). Thus
    
    \[P(m\leq 6) = \int_{-\infty}^6 dP(m) = \frac{3\ln 10}{5}\int_{-\infty}^610^{0.6(m-m_{lim})}dm\]

    Using the integral given in the statement,

    \[P(m\leq 6) = \frac{3\ln 10}{5}2881.6 e^{-1.381 m_{lim}}\]

    Simplifying the numerical factors,

    \[\boxed{P(m\leq 6) \approx 3981.08 e^{-1.381 m_{lim}}} \]

    Out of curiosity, taking \(m_{lim}\approx 40\) as an estimate, we would have

    \[P(m\leq 6)\approx 10^{-20}\]

    Even in the simplest model of the universe, our capability and our insignificance prevail.
    \end{alternativas}

\end{pssolution*}
\end{pproblem}

\pts{4}
\begin{pproblem}
    The \textit{surface brightness}, flux per solid angle per frequency, \(B_\nu\), is given by Planck's law:
    \[B_\nu = \frac{dF}{d\nu d\Omega} = \frac{2h\nu^3}{c^2(e^{\frac{h\nu}{k_BT}}-1)}\] 
    \begin{alternativas}
        \item Find an expression for:
        \[B_\lambda = \frac{dF}{d\lambda d\Omega}\]
        \item Starting from \(B_\lambda\), find the peak wavelength \(\lambda_{max}\) emitted by a star of temperature \(T\) (you will have to solve something numerically).
        \item At low frequencies we have the Rayleigh-Jeans approximation. Obtain an expression for \(B_\nu\) at low frequencies.
        \end{alternativas}

\begin{pssolution*}{}{}
    \begin{alternativas}
        \item We have
        \[B_\lambda = \frac{dF}{d\lambda d\Omega} = \frac{dF}{d\nu d\Omega}\frac{d\nu}{d\lambda}\]

        Using the fundamental wave relation \(c=\nu\lambda \rightarrow \nu = c/\lambda\), we arrive at

        \[B_\lambda = \frac{2h(c/\lambda)^3}{c^2(e^{\frac{hc}{\lambda k_BT}}-1)}\frac{d}{d\lambda}\left(\frac{c}{\lambda}\right) = \frac{2hc^2}{\lambda^5}\frac{1}{(e^{\frac{hc}{\lambda k_BT}}-1)}\]

        \item To solve this part we need the maximum of \(B_\lambda\). Differentiating \(B_\lambda\) directly is extremely tedious. A more efficient approach is to take the natural logarithm of \(B_\lambda\) and differentiate that. At a maximum, both procedures give the same result.
        
        \[\ln B_\lambda = -5\ln \lambda - \ln (e^{hc/\lambda k_BT}-1) + C\]

        where \(C\) is a constant that absorbs the \(\ln\) of the other constants in \(B_\lambda\). Differentiating further,

        \[\frac{d \ln B_\lambda}{d\lambda} = -\frac{5}{\lambda} + \frac{e^{hc/\lambda k_BT}}{e^{hc/\lambda k_BT}-1}\frac{hc}{k_BT\lambda^2}=0\]

        Defining \(x = \frac{hc}{k_BT\lambda}\),

        \[\frac{xe^x}{e^x-1}=5\]

        \[x=5e^{-x}(e^x-1) = 5(1-e^{-x})\]

        Solving by iteration for \(x\) gives \(x\approx 4.965\). Returning to the definition of \(x\),

        \[4.965 = \frac{hc}{k_B T\lambda_{max}}\rightarrow \boxed{\lambda_{max} \approx \frac{2.989 \cdot 10^{-3}}{T}}\]

        This is Wien's law, commonly written \(\lambda = b/T\), where \(b\approx 2.989 \cdot 10^{-3}\).

    \end{alternativas}
\end{pssolution*}
\end{pproblem}

\pts{3}
\begin{pproblem}
    The luminosity of a secondary body depends on the visible illuminated area of the object. In this question we briefly study that phenomenon.
    \begin{alternativas}
        \item  Consider the following situation:
        \begin{figure}[H]
            \centering
            \includegraphics[width=0.8\linewidth]{imagens/fotometria 1.png}
            \caption{Sun-Earth-Moon diagram}
        \end{figure}

        Find an expression for the ratio \(\Phi\) of the illuminated area, as a function of \(\alpha\), to the total area of the planet.
        
        \item Find the angles \(\alpha\) at which the Moon is in the phases new, waxing, full, and waning, respectively.
    \end{alternativas}

\begin{pssolution*}{}{}
    To solve the problem, follow the diagram below:
    \begin{figure}[H]
       \centering
       \includegraphics[width=0.8\linewidth]{imagens/fotometriaplanetaria.png}
       \caption{Source: Introduction to Planetary Photometry} 
    \end{figure}

    On the left is a sketch of how the planet appears in the sky, and on the right a view "from above". The area we see is half the area of the circle plus half the area of an ellipse with semi-major axis \(R\) and semi-minor axis \(R\cos\alpha\). Thus

    \[\Phi = \frac{\pi R^2/2 + \pi R^2\cos\alpha/2}{\pi R^2} \ \boxed{\therefore \Phi = \frac{1+\cos\alpha}{2}}\]

    At new moon, \(\Phi = 0\); at waxing, \(\Phi = 1/2\); at full moon, \(\Phi = 1\); and at waning, \(\Phi = 1/2\). Thus

    \[\boxed{\alpha_{nova} = 180^\circ, \ \alpha_{crescente} = 90^\circ, \ \alpha_{cheia} = 0^\circ, \ \alpha_{minguante} = 90^\circ}\]
\end{pssolution*}
\end{pproblem}

\pts{5}
\begin{pproblem} (Magna notes)
    In this problem we model the effect of the atmosphere on Earth. Suppose that the Sun
is a blackbody of temperature \(T_1\) and radius \(R_1\). Earth is a sphere located
at a distance \(R\) from the Sun, with radius \(R_3\). Earth's emissivity is \(\epsilon_3\).
\begin{alternativas}
    \item If Earth had no atmosphere, determine its equilibrium temperature, \(T_3\).
    \item  Now we include the atmosphere. Model it as a spherical shell of gas,
    with emissivity \(\epsilon_2\) and outer radius \(R_2>R_3\), concentric with Earth. In thermal equilibrium,
    its absorptivity at ultraviolet and infrared wavelengths is \(\epsilon_2\). The
    atmosphere transmits a fraction \(t\) of the ultraviolet radiation but is completely opaque to
    infrared radiation. Assuming that the Sun emits ultraviolet light while Earth emits and re-emits
    in the infrared, determine the temperature \(T_2\) of the atmosphere and the temperature \(T_3\) of Earth in thermodynamic
    equilibrium. Assume that the atmosphere is a perfect thermal conductor, so that
    all radiation incident on it is distributed uniformly over its surface.
\end{alternativas}

\begin{pssolution*}{}{}
    \begin{alternativas}
        \item In thermodynamic equilibrium, the power absorbed equals the power emitted. The emissivity and the absorptivity are also equal. Thus
        \[\frac{4\pi R_1^2\sigma T_1^4}{4\pi R^2}\pi R_3^2\epsilon_3 = 4\pi R_3^2\sigma T_3^4\epsilon_3\]

        Solving for \(T_3\),

        \[\boxed{T_3 = T_1\left(\frac{R_1^2}{4\pi R^2}\right)^{1/4}}\]

         \item The temperature of the atmosphere is
         
         \[T_2 = T_1\left(\frac{R_1^2}{4\pi R^2 }\right)^{1/4}\]

         The reasoning is analogous to the previous part. The luminosity transmitted by Earth is

         \[L_2 = {4\pi R_2^2\sigma T_2^4\epsilon_2} t\]

         Earth reflects all of this radiation, and the atmosphere reflects only a fraction t back, so the total luminosity is

         \[L_{2,T} = 4\pi R_2^2\sigma T_2^4\epsilon_2 (t+t^2+t^3+...)\]

         The expression in parentheses is an infinite geometric series, with first term \(q_1 = t\) and common ratio \(r = t\). Since \(t<1\), the sum is

         \[(t+t^2+t^3+...) = \sum_{n=1}^{\infty}t^n = \frac{t}{1-t}\]

         Equating the two luminosities,

         \[L_{2,T} = L_3 \rightarrow \frac{4\pi R_2^2\sigma T_2^4 \epsilon_2 t}{1-t} = 4\pi R_3^2\sigma T_3^4\]

         Solving for \(T_3\),

         \[T_3 = T_2\left(\frac{R_2^2\epsilon_2}{R_3^2}\frac{t}{1-t}\right)^{1/4}\]

         Substituting \(T_2\),

         \[\boxed{T_3 =T_1\left(\frac{R_1^2R_2^2}{4\pi R^2R_3^2}\frac{\epsilon_2t}{1-t}\right)^{1/4}}\]


    \end{alternativas}
    
\end{pssolution*}
\end{pproblem}


\pts{2}
\begin{pproblem}(Problem set 2 - 2021)
    The Ring Nebula (M57) has an apparent magnitude of 9 and an
angular diameter of 2$^{\prime}$ for an observer on Earth. What would be the apparent magnitude of the night
sky of a planet orbiting a star exactly at the center of M57?
\begin{pssolution*}{}{}
    The night sky covers a solid angle of \(2\pi\) sr, while the solid angle we see is \(\Omega \approx \pi \theta^2\), where \(\theta\) is the angular radius. Since the flux is proportional to \(\Omega\),

    \[m_{ceu}-m = -2.5\log\left(\frac{\Omega_{ceu}}{\Omega}\right)\]

    To convert \(\Omega\) from arcmin\(^2\) to sr, multiply by \(\left(\frac{\pi}{180\cdot60}\right)^2\). Thus
    
    \[m_{ceu} = 9 - \log\left(\frac{2\pi}{\pi(\frac{\pi}{180\cdot60})^2}\right)\]

    \[\boxed{m_{ceu} = -9.43}\]
\end{pssolution*}
\end{pproblem}


\pts{4}
\begin{pproblem} (Adapted from Problem set 4 - 2021)
    The Cherenkov effect was first detected by the Soviet scientist
    Pavel Cherenkov in 1937. Later, together with his colleagues I. E. Tamm and
    I. M. Frank, he gave the phenomenon a physical interpretation and was awarded the Nobel Prize in Physics
    in 1958. Before a mathematical treatment, we first need to understand the principle a little better.

    When high-energy charged particles travel through a dielectric medium, atoms can be excited if
    the particle speed exceeds the phase velocity ($\frac{c}{n}$). Those atoms, in
    turn, emit electromagnetic radiation when they return to the ground state. The emitted waves
    spread spherically and, when superposed, form a cone of opening angle $2\alpha$,
    as the figure below shows.
    \begin{figure}[H]
      \centering
      \includegraphics[width=0.9\linewidth]{imagens/cherenkov.png}  
      \caption{Radiation mechanism of the Cherenkov effect}
    \end{figure}

    This effect is similar to a jet moving at supersonic speed: it follows the same
    principle as the Mach cone, but with light. We now develop the mathematical model of the Cherenkov effect.

    \begin{center}
        \textbf{Part A - Theoretical Model}    
    \end{center}
    
    Consider a particle moving at relativistic speeds through a medium of refractive index
    $n$. Its rest mass is $m_0$, and it has linear momentum $p$ and speed $v$. At a
    given moment, a photon is emitted at an angle $\alpha$, as Figure 1 shows.

    \begin{alternativas}
        \item Letting $\mu$ be the frequency of the emitted photon, determine its linear
        momentum, $p_\mu$, and its energy, $E_\mu$. Your answer must be in terms of $n$, $\mu$, and physical constants.
        \item Find an expression for the linear momentum of the particle after it emits
    the photon, in terms of $p_\mu$, $p$, and $\alpha$.
        \item Letting $\beta_n = \frac{c}{vn}$, prove that the relation below holds:
        \begin{equation}
            \cos\alpha = \frac{1}{\beta_n}
        \end{equation}
        \item Assuming that linear momentum and energy are conserved, determine the
        minimum speed at which the Cherenkov effect can occur.
        \textbf{Hint:} Compared with the other parameters, the factor $(n^2-1)h\mu$ may be neglected.

    \begin{center}
        \textbf{Part B - Nuclear Reactions}    
    \end{center}
    
    The proton-proton chain is a sequence of fusion reactions that convert hydrogen into helium.
    One possible branch of the proton-proton chain is $pp \text{ IV}$, in which, in principle, a
    helium-3 atom reacts directly with a proton, according to the reaction below:

    \begin{equation}
        ^3\text{He} + ^1\text{H} \rightarrow ^4\text{He} + \nu + \_\_
    \end{equation}

        \item  State the nuclear conservation law you use, and identify the missing particle
        in the blank of the reaction above.

        \item  State the nuclear conservation law you use, and identify the missing particle
        in the blank of the reaction below:

        \begin{equation}
            \pi^- \rightarrow \mu^- + \_
        \end{equation}
        
        \textbf{Data:} Pion mass: $140 \ \text{MeV}/c^2$, muon mass: $106 \ \text{MeV}/c^2$.
    \end{alternativas}

\begin{pssolution*}{}{}
    \begin{center}
        \textbf{Part A - Theoretical Model}    
    \end{center}

    \begin{alternativas}
        \item For a photon, the de Broglie relations give
        \[E = h\mu, \ \ p = \frac{h}{\lambda}\]
        where \(h\) is Planck's constant.

        Using the fundamental wave relation, \(\lambda = v/\mu = \frac{c}{n\mu}\), so

        \[\boxed{E = h\mu, \ \ p = \frac{nh\mu}{c}}\]

        \item Since the total momentum is conserved, suppose the particle moves with momentum \(p\) along the \(x\) axis before emitting a photon. Then
        
        \[p = p_\mu\cos\alpha + p'_x\]

        \[p_\mu\sin\alpha = p'_y\]

        \[p' = \sqrt{p_x^2 + p_y^2}\]

        Solving for \(p'\),

        \[p'^2 =(p - p_\mu\cos\alpha)^2 +p_\mu^2\sin^2\alpha \]

        \[p'^2 = p^2 - 2p_\mu p \cos\alpha + p_\mu^2\]

        Finally,

        \[p' = \sqrt{p^2-2pp_\mu\cos\alpha+p^2_\mu}\]

        \item Consider the following figure.
        
        \begin{figure}[H]
            \centering
            \includegraphics[width=0.7\linewidth]{imagens/cherenkov.1.png}
            \caption{Cherenkov diagram}
        \end{figure}

        The radius of the circle is the distance traveled by the photon. The distance from the center of the circle to the end of the cone is \(vt\), so

        \[\boxed{\cos\alpha = \frac{ct}{nvt} = \frac{c}{nv} \equiv \frac{1}{\beta_n}}\]

        \item Setting \(c=1\), energy conservation gives
        \[\sqrt{p^2+m_0^2} = \sqrt{p'^2+m_0^2}+h\mu\]

        Working through this expression,

        \[\sqrt{p'^2 + m_0^2} = \sqrt{p^2+m_0^2} - h\mu\]

        \[p'^2 + m_0^2 = p^2+m_0^2 + h^2\mu^2 - 2h\mu\sqrt{p^2+m_0^2}\]

        Substituting \(p'\),

        \[p^2-2pp_\mu\cos\alpha+p^2_\mu = p^2+h^2\mu^2 - 2h\mu\sqrt{p^2+m_0^2}\]

        Solving for \(\cos\alpha\),

        \[\cos\alpha = \frac{p_\mu^2 - h^2\mu^2+2h\mu\sqrt{p^2+m_0^2}}{2pp\mu}\]

        Simplifying,

        \[\cos\alpha = \frac{p_\mu}{2p} - \frac{h^2\mu^2}{2pp_\mu} + \frac{h\mu\sqrt{p^2+m_0^2}}{pp_\mu}\]
        
        Substituting \(p_\mu = nh\mu\),

        \[\cos\alpha = \frac{nh\mu}{2p} + \frac{h\mu}{2np}+ \frac{\sqrt{p^2+m_0^2}}{np}\]

        \[\cos\alpha = \frac{h\mu(n^2-1)+2\sqrt{p^2+m_0^2}}{2np}\]

        Neglecting the term \(h\mu(n^2-1)\) and substituting \(\cos\alpha\),

        \[\frac{1}{nv} = \frac{\sqrt{p^2+m_0^2}}{np}\]

        Solving for \(v\),

        \[v = \frac{p}{\sqrt{p^2+m_0^2}}\]

        Finally, restoring the factors of \(c\),

        \[\boxed{v = \frac{pc^2}{\sqrt{p^2c^2+m_0^2c^4}}}\]

        \begin{center}
            \textbf{Part B - Nuclear Reactions}    
        \end{center}

        \item Using charge conservation, the left-hand side has \(2+1=3\) protons, but the right-hand side has only \(2\), so the missing particle must contain one proton. The mass number is \(3+1=4\) on the left and \(4\) on the right. The missing particle must therefore have the charge of a proton and a mass much smaller than a proton's. The only particle that fits this description is the \(\boxed{\text{positron, }e^+}\).
            
        \item Lepton number is conserved. The pion is not a lepton, \(n_L=0\), but the muon is a lepton with \(n_L=+1\), so the missing particle must have \(n_L=-1\). The only neutral antilepton associated with the muon and carrying \(n_L=-1\) is the muon antineutrino, \(\bar{\nu}_\mu\). The complete reaction is therefore
        \[\boxed{\pi^{-}\rightarrow \mu^- + \bar{\nu}_\mu}{}\]

    \end{alternativas}

\end{pssolution*}
\end{pproblem}


\pts{3} 
\begin{pproblem}(Problem set 3 - Vinhedo 2022)
Juvelino, observing directly from his observatory in Paris, France, monitors the star Polaris ($\alpha\text{UMi}$). His goal is to find the color temperature $T_c$ of the star. Some of the data he has on his target are:

\begin{itemize}
    \item Apparent magnitude in the $V$ band: $V = 1.98$;
    \item Absolute magnitude in the $V$ band: $M_V = -3.60$;
    \item Absolute magnitude in the $B$ band: $M_B = -3.19$.
\end{itemize}

Using the information provided, help Juvelino!

\begin{alternativas}
    \item From many observations, Juvelino determined that the interstellar extinction in the $V$ band toward Polaris is $a_V = 5.8 \, \text{mag/kpc}$. Find the distance, in pc, from $\alpha\text{UMi}$ to Earth.
    \item Using the empirical relation
    \begin{equation}
        \frac{A_V}{E_{B-V}} = 3.0
    \end{equation}
    where $A_V$ is the total interstellar extinction in the $V$ band and $E_{B-V}$ is the $B-V$ color excess, find the $B-V$ color index of the observed star.
    \item Show the relation
    \begin{equation}
        T_c = \frac{7009}{(B-V) + 0.47}
    \end{equation}
    in which the color temperature is given in kelvin. To do so, use the fact that stars of spectral class $A0$ have $(B-V) = 0$ and $T_c = 15000 \, \text{K}$. Also use the wavelengths of the $B$ and $V$ bands, $\lambda_B = 440 \, \text{nm}$ and $\lambda_V = 548 \, \text{nm}$, respectively. Justify any approximations you make.

    \textbf{HINT: Planck's law may be useful}
    \item Find the color temperature of Polaris.
\end{alternativas}
\begin{pssolution*}{}{}
    \begin{alternativas}
        \item The expression that correctly relates the magnitudes is
        \[V - M_V = 5 \log d - 5 + a_V d\]

        This equation can be solved only by iteration:

        \[d = \frac{V-M_V - 5\log d + 5}{a_v}\]

        Iterating, we obtain \(\boxed{d\approx 100\text{ pc}}\).

        \item By definition, \(E_{B-V} = A_B - A_V\). Thus
        
        \[\frac{A_V}{A_B-A_V} = 3.00\]

        \[A_B = \frac{A_V}{3}+A_V = 0.77\]

        Thus

        \[B-M_B = 5\log d - 5 + A_B\rightarrow B = 2.58\]

        By definition, the \(B-V\) index is \(\boxed{U_{B-V} = B-V = 0.60}\)

        \item Using Planck's law,
        
        \[\frac{B_B}{B_V} = \frac{F_B}{F_V} \left(\frac{\lambda_V}{\lambda_B}\right)^5 \frac{e^{\frac{\beta hc}{\lambda_V}-1}}{e^{\frac{\beta hc}{\lambda_B}-1}}\]

        Using the Wien approximation,

        \[\frac{F_B}{F_V} = \left(\frac{\lambda_V}{\lambda_B}\right)^5 \frac{e^{\frac{\beta hc}{\lambda_V}}}{e^{\frac{\beta hc}{\lambda_B}}}\]

        Using Pogson's equation to compare the magnitudes,

        \[B-V = -2.5\log\left(\frac{F_B}{F_V}\right)+C\]

        A constant is included because the two bands have different wavelengths. Working through the expression,

        \[B-V = -2.5\log\left(\left(\frac{\lambda_V}{\lambda_B}\right)^5 \frac{e^{\frac{\beta hc}{\lambda_V}}}{e^{\frac{\beta hc}{\lambda_B}}}\right)+C\]

        Working through it further,

        \[B-V = -2.5\log\left(\frac{\lambda_V}{\lambda_B}\right)^5 + 2.5 \beta hc\left(\frac{1}{\lambda_B}-\frac{1}{\lambda_V}\right)\log e + C\]

        Using \(\beta = \frac{1}{k_B T_c}\) and substituting the values,

        \[B-V = -1.19 + \frac{7009}{T_c}+C\]

        For type \(A_0\), \((B_V)=0\) and \(T_c = 15000\)K. Thus

        \[-1.19 + \frac{7009}{15000} + C = 0 \rightarrow C = 0.72\]

        Thus

        \[B-V = -0.47 + \frac{7009}{T_c} \rightarrow \ \ \boxed{T_c = \frac{7009}{(B-V) + 0.47}}\]

        \item From the previous equation,
        
        \[\boxed{T_c = \frac{7009}{0.60 + 0.47} \approx 6600 \text{ K}}\]
    \end{alternativas}    

\end{pssolution*}
\end{pproblem}


\pts{4}
\begin{pproblem}
    (Problem set 4 - 2020) Frightening as it is, Juvelino's banana plantation has a perfect sky for one of his two passions: astrophotography. He has a telescope 200 mm in diameter with a super CCD attached. Its parameters are: pixel side \(5\ \mu\text{m}\), plate scale \(3.2\ \text{arcmin/mm}\), quantum efficiency 97\%, RON (read-out noise) 1 count/pixel, and DN (thermal noise) 1 count/(pixel.hour). He prefers to observe in the \(V\) band (central wavelength \(\lambda_V = 5500\ \text{Å}\), bandwidth \(\Delta \lambda = 820\ \text{Å}\)), because he remembers a magic number associated with it: the photon flux of a magnitude-0 object in this band is \(\phi_0 = 1000\ \text{photons/}(\AA\cdot \text{cm}^2)\). Juvelino has two favorite targets: a star of magnitude \(m_V = 4.51\) and a globular cluster of uniform surface brightness in \(V\) of \(19.5\ \text{mag/arcsec}^2\) and angular diameter \(\theta = 40''\).

    \begin{alternativas}
        \item The night-sky brightness at Juvelino's farm is \(21.9\) mag/arcsec\(^2\) in the V band, and the astronomical seeing is \(\theta_0=1.2''\). Calculate the signal-to-noise ratio obtained in a photo of his favorite star with exposure time \(t_{exp}=30\)s.
        \item Calculate the zero-point magnitude of Juvelino's telescope + CCD system, using the data for the star.
        
        Juvelino's other great passion is Jiang, a Chinese university student who lives in Beijing. After years apart, Jiang suggests that Juvelino move to China. Before deciding, he looks up the night-sky brightness in the V band in his beloved's city, which is \(19.3\) mag/arcsec\(^2\). That will affect imaging of the cluster, which is also visible from Jiang's home.

        \item Juvelino cannot live without photographing the cluster. He has inviolable rules for imaging this object, wherever he does it: reach a signal-to-noise ratio $S/N=10$, and take the same number $X$ of photos of the cluster in one night. $X$ is the number of photos allowed by the sky conditions at his banana plantation. That number of frames per night is limited by the CCD battery (which lasts $3.2$ hours) and by the exposure time required for each photo. Help him with this dilemma! Do the calculations, find how many photos of the cluster he can take in one night (in Beijing and at the farm), and conclude: will he see his beloved again, or will this love be too noisy?
    \end{alternativas}
\begin{pssolution*}{}{}
    The official solution can be found \href{https://drive.google.com/file/d/1q81Ba0RXA_L845xAE0uhVMecuhXAr2Ad/view?usp=sharing}{here}.
\end{pssolution*}
\end{pproblem}


\end{document}
